\section{Construction}\label{sec:txidconstruction}

We have a ledger that processes user transactions.  We also have two one way
functions, which we call $\G$ and $\G'$, and two PRFs, $\FE$ and $\FT$.  There
is an upper bound of the number of transactions a user can submit during any
given epoch, which we call $\Max$.

The protocol operates as follows. Each user is provisioned with a long term
signing key and a long term tracing key. These two keys are cryptographically
bound together via a credential issued by the registration authority.  For each
epoch, the user derives an epoch tracing key as a function of their tracing key
and the epoch identifier. All transactions submitted by the user during that
epoch are tagged with a PRF evaluated with the user's epoch tracing key and an
index in the range $[0,\Max)$.  Each transaction is accompanied by a proof of
the correctness of the tag generation.

For transaction identification, the user's epoch tracing key is disclosed to
the auditor in question.  The auditor can then locally compute all valid
transaction tags corresponding to that user and epoch.

When the system transitions to a new epoch, users derive a fresh epoch tracing
key. Tracing capabilities granted for prior epochs automatically expire.

The protocol uses signatures of knowledge where a proof is also used to
authenticate a message. For proofs used only to certify correctness, we use non
interactive zero knowledge proofs.

\paragraph{Setup.} During setup, the registration authority generates a key
pair for a digital signature scheme. Their public key is communicated to the
ledger operator. The ledger operator generates the public parameters for a
commitment scheme and for four proof systems. The first relation allows users
to prove knowledge of their long term signing and tracing keys, which they do
to register in the system. The second captures correctness of epoch tracing key
generation.  The third allows users to prove possession of a valid registration
credential, and consistency between the tracing key certified by the
registration credential and the tracing key contained in the commitment.  The
fourth relation enables users to prove that a given epoch tracing key and
transaction tag were correctly derived from the tracing key contained in the
same commitment.

The registration authority generates $(\rask, \rapk) \leftarrow
\DSgen(\secparam)$ and communicates $\pk$ to the ledger operator.  The ledger
operator calls $ \ck \leftarrow \ComGen(\secparam)$ and produces four sets of
public parameters:

\begin{itemize}

  \item First, we have $ \inst_0 = (\G, \G', \PSK, \PTK)$, and $\witness_0 =
    (\sk, \utk)$ for the relation defined as follows:
   %
    \begin{equation}  \label{eq:rel0} \PSK = \G(\sk) \quad \wedge \quad \PTK =
    \G'(\utk).  \end{equation}

  \item Next, we have $\inst_1 = (\FE, \G', \PTK, \epoch,  \ek)$, and
    $\witness_1 =  \utk$, for the following relation:
    %
    \begin{equation} \label{eq:rel1} \PTK = \G'(\utk) \quad \wedge \quad \ek =
    \FE(\utk, \epoch).  \end{equation}

  \item The third are defined $\inst_2 = (\G, \G', \rapk, \ck, \com)$, and
    $\witness_2 = (\signature, \sk, \utk, \openinfo)$, for relation given as
    follows:
  %
    \begin{equation} \label{eq:rel2} \DSverify(\rapk, \G(\sk) \| \G'(\utk),
    \signature) = 1 \ \wedge  \ \ComOpen(\ck, \openinfo, \com) = \utk.
  \end{equation}

\item  And finally, we have $\inst_3 = (\FE, \FT, \ck, \com, \epoch, \Max,
  \Tag)$, and $\witness_3 =  (\utk, \openinfo, \ek, \gamma)$, for proving that
  the following are satisfied:
%
  \begin{gather} \label{eq:rel3} \ComOpen(\ck, \openinfo, \com) = \utk \quad
  \wedge \quad \ek = \FE(\utk, \epoch) \\ \wedge \quad \Tag = \FT(\ek, \gamma)
\quad  \wedge  \quad 0 \leq \gamma < \Max.  \end{gather}

\end{itemize}

The ledger operator then publishes the following:
%
\begin{equation*} \ledger[\crs] \leftarrow (\G, \G', \FE, \FT, \rapk, \ck,
\crs_0, \crs_1, \crs_2, \crs_3, \Max).  \end{equation*}

\paragraph{Registration.}  During registration, the user proves knowledge of
their long term signing and tracing keys. The registration authority then
certifies the corresponding public values, binding them together for future
use.

The registration authority first checks if the user has registered before.  The
user generates signing key pair $(\pk, \sk)$ and tag key $\utk$, computes $\PSK
= \G(\sk)$ and $\PTK = \G'(\utk)$. They then engage in an interactive protocol
with the registration authority.  During this interaction, the user receives a
nonce $\rand$, computes and sends back $\soksig \leftarrow \soksign(\crs_0,
\inst_0, \witness_0, \rand)$, for $(\inst_0, \witness_0)$ as defined above.

If the registration authority accepts $\soksig$, then it calls $\DSsign(\rask,
\PSK \| \PTK)$, records that the user registered with pair $(\PSK, \PTK)$, and
sends the resulting signature $\signature$ to the user. The user's public key
is then defined as $\PSK = \G(\sk)$.

\paragraph{Authorisation.} Before the start of epoch $\epoch$, auditors issue
requests for tracing authorisations to users.  In response, the relevant user
computes their epoch tracing key $\ek = \FE(\utk, \epoch)$, and generates a
zero knowledge proof $\proof \leftarrow \NIZKprove(\crs_1, \inst_1,
\witness_1)$, with $\inst_1$ and $\witness_1$ as defined above. The proof that
the epoch tracing key is correctly derived can be reused throughout an epoch.
To authorise a particular auditor, the user additionally authenticates that
proof using a signature of knowledge.

An authorisation request identifies the auditor and contains pair $(\epoch,
\rand)$. The user grants the auditor the tracing authorisation by sending tuple
$(\PSK, \PTK, \signature, \ek, \proof, \soksig)$, with $\soksig \leftarrow
\soksign(\crs_0, \inst_0, \witness_0, \proof \| \rand)$, and $(\inst_0,
\witness_0)$ as defined above.  If $1 \leftarrow \DSverify(\rapk, \PSK \| \PTK,
\signature)$, and $\proof$ and $\soksig$ are valid, then the auditor accepts
tracing key $\ek$.

\paragraph{Transaction submission.} During transaction submission, the user
proves that the transaction tag was correctly generated from their tracing key
while also authenticating the transaction itself.

The user first prepares confidential payload of their transaction denoted
$\payload$, selects a fresh counter $0 \leq \gamma < \Max$, and computes $\Tag=
\FT(\ek, \gamma)$.  Next, the user computes commitment $\com \leftarrow
\ComCommit(\ck, \utk, \openinfo)$ and generates $\proof \leftarrow
\NIZKprove(\crs_3, \inst_3, \witness_3)$, for $(\inst_3, \witness_3)$ as
defined above. The user produces $\soksig \leftarrow \soksign(\crs_2, \inst_2,
\witness_2, \msg)$, with $\msg = \payload \| \Tag\| \com \| \proof$ and
$(\inst_2, \witness_2)$ as defined above.  The user submits the transaction
$(\payload, \Tag, \com, \proof, \soksig)$ to the ledger for validation.

Once the transaction is appended to the ledger, the transaction validators
check if $\Tag$ has never accompanied a previous transaction that was marked
valid; that $\proof$ is a valid proof for relation $\rel_3$ and current epoch
$\epoch$; and that $\soksig$ is a valid proof on message $\msg = \payload \|
\Tag\| \com \| \proof$ relative to relation~$\rel_2$.

If these checks succeed, the transaction validators mark the transaction as
valid by adding entry $\ledger[\Tag] \leftarrow \payload$.

\paragraph{Transaction identification.} Given key $\ek$ for epoch $\epoch$, the
auditor traces the transactions that the user submitted during that epoch as
follows.  $\forall \ 0 \leq \gamma < \Max:$  the auditor computes $\Tag=
\FT(\ek, \gamma)$ and retrieves entry $\ledger[\Tag]$ if not empty.

\paragraph{Increment Epoch.} The system participants move to epoch $\epoch +
1$.

\begin{theorem}\label{thm:bb} Assuming that $\DS$ satisfies existential
unforgeability, $\FE$ and $\FT$ are psuedorandom, $\G$ and $\G'$ are one way,
and $\SOK$ and the proof system are secure, the protocol described above securely realises
$\F{txid}$.  \end{theorem}
